%------------------------
% Resume Template
% Author : Anubhav Singh
% Github : https://github.com/xprilion
% License : MIT
%------------------------

\documentclass[a4paper,20pt]{article}

\usepackage{latexsym}
\usepackage[empty]{fullpage}
\usepackage{titlesec}
\usepackage[usenames,dvipsnames]{color}
\usepackage{verbatim}
\usepackage{enumitem}
\usepackage[pdftex]{hyperref}
\usepackage{fancyhdr}

\pagestyle{fancy}
\fancyhf{} % clear all header and footer fields
\fancyfoot{}
\renewcommand{\headrulewidth}{0pt}
\renewcommand{\footrulewidth}{0pt}

% Adjust margins
\addtolength{\oddsidemargin}{-0.530in}
\addtolength{\evensidemargin}{-0.375in}
\addtolength{\textwidth}{1in}
\addtolength{\topmargin}{-.45in}
\addtolength{\textheight}{1in}

\urlstyle{rm}

\raggedbottom
\raggedright
\setlength{\tabcolsep}{0in}

% Sections formatting
\titleformat{\section}{
  \vspace{-10pt}\scshape\raggedright\large
}{}{0em}{}[\color{black}\titlerule \vspace{-6pt}]

%-----------------------------
%%%%%%  CV STARTS HERE  %%%%%%

\begin{document}

%----------HEADING-----------------
\begin{center}
  \textbf{{\Large Khimraj Suthar}} \\
  Bengaluru, India \hspace{1em}
  \href{mailto:khimrajsuthar@gmail.com}{khimrajsuthar@gmail.com} \hspace{1em}
  \href{tel:+918290285519}{+91-8290285519} \\
  \href{https://www.linkedin.com/in/khimraj-suthar-20b7aa122/}{linkedin.com/in/khimraj-suthar-20b7aa122} \hspace{1em}
  \href{https://github.com/khimraj}{github.com/khimraj} \hspace{1em}
  \href{https://khimrajsuthar.com}{khimrajsuthar.com}
\end{center}

%-----------SUMMARY-----------------
\section{Summary}
GenAI Engineer with 5+ years building production AI systems spanning agentic workflows, retrieval systems, model evaluation, fine-tuning, and AI infrastructure across OpenAI, Claude, Gemini, and Llama ecosystems. Delivered enterprise-scale solutions processing 4M+ records at 100K/day throughput, serving high-throughput workloads, and driving organization-wide AI adoption.

%-----------SKILL SUMMARY-----------------
\vspace{1pt}
\section{Skills}
	\noindent\textbf{LLM Systems:} RAG, Multi-Agent Systems, Agent Orchestration, A2A, MCP, Function Calling, Prompt \& Context Engineering, Fine-tuning, Model Evaluation, LLM Routing, Guardrails, Observability \\
	\textbf{Frameworks \& Models:} LangGraph, LangChain, CrewAI, ADK, Glean, Vertex AI, Claude Code $\cdot$ GPT, Claude, Gemini, Llama \\
	\textbf{Infrastructure \& Data:} AWS, GCP, Azure, Kubernetes, Docker, Airflow, Arize, CI/CD $\cdot$ PostgreSQL, Redis, Qdrant, Vector DB \\
	\textbf{Languages \& ML:} Python, C++, SQL $\cdot$ TensorFlow, Scikit-Learn, XGBoost, Deep Learning, NLP, Computer Vision
 
%-----------EXPERIENCE-----------------
\section{Experience}
  \noindent\textbf{GenAI Engineer} \hfill \textit{December 2024 - Present}\\
    Docusign \hfill Bengaluru, India
    \vspace{-2pt}
    \begin{itemize}[label=$\bullet$]
        \item Designed a multi-stage LLM classification pipeline with model routing, confidence scoring, human-review workflows, and offline evaluation to track accuracy, processing 4M+ records at 100K/day throughput.
        \vspace{-4pt}
        \item Built shared GenAI components including LLM routing with cost and latency optimization, MCP services, agent identity, RAG pipelines, and observability integrations that accelerated development across AI applications.
        \vspace{-4pt}
        \item Led development and rollout of Glean-based productivity agents, including automated daily action items, adopted by 50\%+ of employees.
        \vspace{-4pt}
        \item Developed a Text-to-SQL pipeline using LangChain and vector databases, achieving 94\% query accuracy and reducing reporting turnaround by 80\%.
        \vspace{-4pt}
        \item Architected personalized email generation agents using CrewAI and LangChain, delivering 10K+ tailored emails and increasing click-through rates by 14\%.
      \end{itemize}
 
  \noindent\textbf{Software Engineer, LLM} \hfill \textit{December 2023 - November 2024}\\
    Webb.ai \hfill Remote, India
    \vspace{-2pt}
    \begin{itemize}[label=$\bullet$]
        \item Architected multi-agent and RAG systems for automated DevOps/SRE troubleshooting in cloud environments, handling 100 parallel requests per second in production.
        \vspace{-4pt}
        \item Fine-tuned LLMs using LoRA on tool-call datasets, built evaluation pipelines to benchmark function-call accuracy on held-out benchmarks, improving it from 78\% to 89\%.
        \vspace{-4pt}
        \item Drove a 60\% reduction in DevOps/SRE debugging time and 50\% cost savings through LLM-driven root cause analysis and incident resolution workflows.
      \end{itemize}
 
    \noindent\textbf{Software Engineer, Machine Learning} \hfill \textit{June 2021 - November 2023}\\
    Ivy Homes \hfill Bengaluru, India
    \vspace{-2pt}
    \begin{itemize}[label=$\bullet$]
        \item Led development of OCR, Entity Extraction, Document Classification, Text Clustering, and Pricing/Valuation models, reducing property valuation error from 15\% to 8\% MdAPE.
        \vspace{-4pt}
        \item Built an LLM-powered real estate chatbot that automated 90\% of off-hours user queries. Developed data pipelines processing 3M+ records and secured internal dashboards and business data through a self-hosted VPN layer.
      \end{itemize}
 
    \noindent\textbf{Software Development Engineer Intern} \hfill \textit{May 2020 -- July 2020}\\
    Amazon \hfill Remote, India
    \vspace{-2pt}
		\begin{itemize}[label=$\bullet$]
        \item Built and deployed a lender runbook website for Amazon Pay Later on AWS, reducing partner prep time by 80\%, won a 3-day hackathon with 30+ participants and received a full-time offer.
		\end{itemize}
 
%-----------EDUCATION-----------------
\vspace{1pt}
\section{Education}
\noindent\textbf{Bachelor of Technology - Computer Science and Engineering} \hfill \textit{August 2017 - June 2021}\\
Malaviya National Institute of Technology (MNIT) \hfill Jaipur, India
 
%-----------PUBLICATIONS-----------------
\vspace{1pt}
\section*{Publications}
\begin{itemize}[leftmargin=*]
    \item Khimraj, P. K. Shukla, A. Vijayvargiya, R. Kumar. \href{https://ieeexplore.ieee.org/document/9117456}{\textbf{Human Activity Recognition using Accelerometer and Gyroscope Data from Smartphones}}. \textit{ICONC3}, 2020.
    \vspace{-4pt}
    \item Vijayvargiya A, Khimraj, Kumar R, Dey N. \href{https://link.springer.com/article/10.1007/s13246-021-01071-6}{\textbf{Voting-based 1D CNN model for human lower limb activity recognition using sEMG signal}}. \textit{Phys Eng Sci Med}, 2021.
\end{itemize}
 
\end{document}